\documentclass{article}
\usepackage{graphicx} % Required for inserting images

\title{Freelancing}
\author{Riccardo Parrino}
\date{May 2025}

\begin{document}

\maketitle

\section{Modello di business}

\section{Comunicazione}

\subsection{Comunicazione interna}
Rivolta ai dipendenti e ai collaboratori. Ha come obiettivi aumentare la motivazione, la condivisione della cultura aziendale, il coordinamento interno e la gestione del cambiamento. Viene fatto attraverso newsletter, intranet, riunioni, formazione ed eventi interni.

\subsection{Comunicazione esterna}
Rivolta al pubblico esterno all'azienda. Ha come obiettivi promuovere l'immagine aziendale, costruire relazioni, migliorare la reputazione.

Si articola in comunicazione istituzionale, comunicazione commerciale/di marketing, comunicazinoe finanziaria e comunicazione di crisi.

La comunicazione istituzionale cura l'indentit\'a e l'immagine dell'azienda,  ha come target istituzione, investitori ed opinione pubblica. Usa strumenti com e le relazioni pubbliche, gli eventi, i comunicati stampa ed i rapporti di sostenibilit\'a.

Per comunicazione istituzionale, pi\'u precisamente, si intende la comunicazione realizzata in moda organizzato da un'istituzione o dai suoi rappresentanti e diretta alle persone e ai gruppi dell'ambiente sociale in cui svolge la sua attivit\'a. Ha come obiettivo stabilire relazioni di qualit\'a tra l'isituzione e il pubblico con cui si relaziona, per conseguire notoriet\'a sociale e immagine pubblica adeguate ai fini e alle attivit\'a dell'istituzione stessa (Comunicazione pubblica). Nell'ambito pubblicitario, il termine indica la comunicazione volta a promuovere l'immagine di un organismo, commerciale e non, inteso nel suo insieme (corporate image).

La comunicazione commerciale/di marketing mira a promuovere i prodotti o servizi ha come target i clienti attuali e potenziali e usa strumenti come pubblicit\'a , promozioni, direct marketing, digital marketing e social media.

La comunicazione finanziaria \'e rivolta agli investitori, analisti e azionisti, ha contenuti come i risultati economici, bilanci e operazioni finanziarie. 

Infine la comunicazione di crisi mira a comunicaer le situazioni critiche che possono danneggiare l'immagine o la reputazione dell'impresa ed ha come obiettivi la trasparenza, il contenimento dei danni ed il recupero della fiducia.

\subsection{Marketing}
Una visione di insieme del marketing:

\begin{itemize}
    \item Markeing planning and strategy, Business mission and strategy, marketing audit
    \item Marketing objective, strategy, actions
    \item Marketing Evaluation of performance, business performance, metrics
    \item The scope of marketing
\end{itemize}

Understanding Customer Behaviour:
\begin{itemize}
    \item The dimensions of customer behaviour: Who buys, how they buy, what are the choice criteria
    \item  Influences on consumer behaviour, on organizational buying behaviour
\end{itemize}

Marketing Research and Customer Insights
\begin{itemize}
    \item Market research
    \item The role of customer insight
    \item Marketing information systems, market intelligence
\end{itemize}

Market Segmentation, Targeting and Positioning
\begin{itemize}
    \item Segmenting consumer markets, criteria
    \item Segmenting organizational markets, criteria
    \item Target marketing, Positioning, Repositioning
\end{itemize}

Creating Customer Value
\begin{itemize}
    \item Value through Products and Brands
    \item Value through Services, Relationships and Experiences
    \item Value through Pricing
\end{itemize}

Delivering and Managing Customer Value
\begin{itemize}
    \item Delivering Customer Value
    \item Mass Communications Techniques
    \item Direct communications techniques
    \item Digital Marketing
\end{itemize}

\end{document}